% TOP_PROBLEM: {"id": 3023, "title": "Kirby Problem 5.16", "category_id": 7, "category": {"id": 7, "name": "topology", "display_name": "Topology", "description": "Properties preserved under continuous deformations.", "slug": "topology", "order_index": 7, "created_at": "2026-07-31T15:26:25.671Z"}, "status": "open"}
% FIRST_POSTED: null
% ENTRY_KIND: statement_only
% Imported from: batch14.tex; UnsolvedMath ID 3023
\documentclass[11pt]{article}
\usepackage[margin=25mm]{geometry}
\usepackage{fontspec}
\setmainfont{Latin Modern Roman}
\usepackage{amsmath,amssymb,mathtools}
\usepackage{unicode-math}
\setmathfont{Latin Modern Math}
\usepackage{xurl,hyperref}
\hypersetup{colorlinks=true,urlcolor=blue}
\setcounter{secnumdepth}{0}
\setlength{\parindent}{0pt}
\setlength{\parskip}{5pt}
\setlength{\emergencystretch}{3em}
\newcommand{\sref}[2]{\hyperlink{source:#1:#2}{[S#2]}}
\newcommand{\cataloguescope}[1]{\paragraph{Recorded target.}#1}
\begin{document}
% UnsolvedMath ID 3023; source catalogue code KP-5.16
\setcounter{section}{3022}
\section{Kirby Problem 5.16}\label{3023}
{\small\sffamily UnsolvedMath ID 3023 · Topology}
\subsection{Definitions and mathematical statement}

\paragraph{Shared notation.}$\mathbb N=\{1,2,\ldots\}$, and $\mathbb Z,\mathbb Q,\mathbb R,\mathbb C$ denote the integers, rationals, reals and complex numbers. Any other conventions are specified locally.


Fix $R=\mathbb Z$ or a field, with an effective coefficient presentation when decidability is interpreted in the Turing model. A differential graded $R$-algebra (DGA) is a unital algebra $A=\bigoplus_{j\in\mathbb Z}A_j$ with an $R$-linear differential $d$ of degree $-1$, $d^2=0$, and $d(ab)=d(a)b+(-1)^{\deg a}a\,d(b)$ for homogeneous $a,b$. Inputs specify finitely many homogeneous generators (in arbitrary integer degrees) and the finite polynomial or word expressions for their differentials. The underlying graded algebra is either free associative or free graded-commutative over $R$; in the latter case $ab=(-1)^{\deg a\deg b}ba$.

An elementary change of free generators fixes all but one and replaces $x_i$ by $u x_i+P$, where $u\in R^\times$ and $P$ is homogeneous of the same degree in the other generators. A tame isomorphism is a composition of these changes and a relabeling, commuting with the differentials. Stabilization adjoins generators $e,f$ with $\deg e=\deg f+1$, $d(e)=f$, $d(f)=0$. Stable tame isomorphism means tame isomorphism after finitely many stabilizations. A DGA map is a quasi-isomorphism if it induces an isomorphism on homology. In the source's convention, $A$ and $B$ are quasi-isomorphic if there is an $R$-DGA $C$ and quasi-isomorphisms $C\to A$ and $C\to B$; no finite-generation condition on $C$ is imposed. \sref{3023}{2} Derived Morita equivalence means equivalence of the derived categories of differential modules as triangulated categories. \sref{3023}{3}

Over fields not supplied with effective encodings, a model of coefficient computation must be specified; the quoted paper distinguishes Turing computation from real/complex algebraic computation.
\paragraph{Statement recorded in the supplied source.}
For each of the two classes of underlying graded algebras and each indicated coefficient ring, are there algorithms deciding, from finite descriptions of $A,B$, whether:
\begin{enumerate}
\item[(a)] $A$ and $B$ are stable tame isomorphic;
\item[(b)] $A$ and $B$ are quasi-isomorphic;
\item[(c)] $A$ and $B$ are derived Morita equivalent?
\end{enumerate}
The numbers and degrees of the generators of $A$ and $B$ may differ, and negative degrees are allowed. \sref{3023}{2}
\subsection{Short English statement}
Determine whether three kinds of equivalence of finitely generated differential graded algebras can be decided, in both the noncommutative and graded-commutative settings.
\subsection{Sources}
\begin{description}
\item[\hypertarget{source:3023:1}{S1}] ULAM.AI and the UnsolvedMath Contributors, \emph{UnsolvedMath: A Curated Collection of Open Mathematics Problems}, record 3023 (KP-5.16). CC BY 4.0. \url{https://huggingface.co/datasets/ulamai/UnsolvedMath}.
\item[\hypertarget{source:3023:2}{S2}] R. İnanç Baykur, Robion C. Kirby and Daniel Ruberman (eds.), \emph{K3: A New Problem List in Low-Dimensional Topology}, Mathematical Surveys and Monographs 295, American Mathematical Society (2026), Problem 5.16, author preliminary version. \url{https://math.berkeley.edu/sites/default/files/surv-295-ruberman-watermarked-author-pdf.pdf}.
\item[\hypertarget{source:3023:3}{S3}] Ciprian Manolescu and Nick Rozenblyum, \emph{Undecidability problems for semifree DG algebras}, arXiv:2605.08122v1 (2026), Theorem 1, its coefficient-ring footnote, and §2. Theorem 1 answers all three noncommutative questions negatively for computable coefficient rings; the paper explicitly leaves the graded-commutative alternative open. \url{https://arxiv.org/abs/2605.08122}.
\end{description}
\label{end:3023}
\end{document}
